% =============================================================================
% Cover Letter for the NL2Semantic2SQL submission.
% Generic single-page format suitable for IJGIS / GeoInformatica /
% Transactions in GIS / ISPRS IJGI; please replace journal name, editor name,
% and submission date as appropriate.
% =============================================================================
\documentclass[11pt]{letter}
\usepackage[margin=1in]{geometry}
\usepackage{times}
\usepackage{url}
\usepackage[hidelinks]{hyperref}

\signature{Anonymous Author(s) \\ \textit{Affiliation withheld for review}}
\address{Anonymous Author(s) \\ \textit{Affiliation withheld for review} \\ \texttt{anonymous@example.com}}

\begin{document}

\begin{letter}{The Editor-in-Chief \\ \emph{[Journal Name]}}

\opening{Dear Editor,}

We are pleased to submit our manuscript, ``\textbf{NL2Semantic2SQL: A Cross-Domain Framework for Natural Language to SQL over Geospatial and Enterprise Data Warehouses},'' for consideration as a regular research article in \emph{[Journal Name]}.

The work addresses a gap in the natural-language-to-SQL literature: existing systems and benchmarks (Spider, BIRD, Spider~2.0, BEAVER) almost exclusively target conventional relational databases, and there is no execution-based benchmark for spatial SQL generation, no systematic study of cross-domain transfer between GIS and warehouse query semantics, and no treatment of cross-lingual querying that combines language and operator heterogeneity. Our manuscript closes these three gaps in one study.

Specifically, the manuscript makes the following contributions:

\begin{enumerate}
\item We propose a unified semantic-to-SQL pipeline that decomposes the task into semantic-layer grounding, schema-aware context construction, and execution-time self-correction. The same pipeline serves both a GIS database (with PostGIS) and a warehouse database derived from BIRD mini\_dev, with domain-specific evidence injected at the grounding stage rather than encoded into the model.
\item We construct a 20-question GIS-oriented benchmark with execution-based evaluation, structured spatial-operator coverage, and a Robustness category that probes safety/refusal behavior. Combined with a 50-question BIRD subset, this gives a dual-track protocol that has not previously been reported.
\item Through three pipeline configurations (baseline LLM, full pipeline with prompt refinements, and full pipeline with MetricFlow-style fact/dimension modeling) we identify the bottleneck on the warehouse side and show that warehouse-specific structural metadata alone closes the gap to baseline parity. A preliminary cross-lingual experiment in which BIRD questions are translated into Chinese further shows that multilingual alias registration is sufficient to support cross-lingual querying with only a four-percentage-point loss.
\end{enumerate}

We believe the paper is a good fit for \emph{[Journal Name]}: it spans GIS-centric methods (geometry-aware grounding, spatial operators, coordinate-system rules), GIS-relevant evaluation (PostGIS, KNN, ST\_Intersection, ST\_Distance), and the broader text-to-SQL methodological discussion.

\medskip
\noindent\textbf{Status of experimental data.} The current draft reports execution accuracy on a 20-question GIS benchmark and a 50-question subset of BIRD mini\_dev. To strengthen statistical robustness, we are currently running the full 500-question version of BIRD mini\_dev on the same pipeline configurations; partial baseline results from this large-scale run are already available, and the complete 500-question numbers will be incorporated into the manuscript before the first round of review concludes. The conclusions of the paper do not depend on the exact 50-vs-500 sample size, only on the relative ordering of the three configurations, which is consistent in our partial results.

\medskip
\noindent\textbf{Reproducibility.} All benchmark questions, gold SQL, prompt templates, semantic-model registrations (including the MetricFlow warehouse models), and per-question execution logs are released with the submission. Each table in the manuscript corresponds to a specific run directory under \texttt{nl2sql\_eval\_results/} that contains the gold SQL, the predicted SQL, the execution comparison, and the per-question token usage.

\medskip
\noindent\textbf{Originality and conflicts of interest.} The manuscript is original work and has not been published or submitted elsewhere. No portion has appeared in any conference or workshop. The authors declare no competing interests.

We thank you for considering our submission and look forward to the reviewers' feedback.

\closing{Sincerely,}

\end{letter}

\end{document}
